Symmetry is one of the most powerful priors we can give a learning model: it marks the variations in the input that should leave the answer unchanged, so the model need not learn them from data. The prior is especially valuable in quantum machine learning, where circuits are small and easily overwhelmed, and recent work shows that building a known symmetry into a variational quantum circuit improves generalization. We ask what comes next: once a model is symmetric, is the only choice how much symmetry to impose, or also which symmetry-respecting interactions it may learn? Treating symmetry as a design space, we relax it toward smaller subgroups and find that weaker constraints recover almost all of the benefit, then keep the symmetry but enrich it with trainable interactions acting directly on the patterns that determine the label. This second move lifts performance well beyond the original symmetric circuit, and a control that scrambles the parameter sharing at matched size confirms the gain comes from the task's geometry, not a smaller model. We demonstrate this on Tic-Tac-Toe board classification, chosen because every relevant symmetry, label, and decisive pattern is known exactly, so the effect can be isolated rather than inferred. A known symmetry is not merely a way to tie parameters; it is a language for deciding which interactions a model may learn, and so it is not the end of architecture design but its beginning.
